% ============================================================================
% LANGENTWURF LE-012: Suchmaschinen
% Profil: S (Senioren 65+) | Tier: 2 (Standard)
% ============================================================================

\documentclass[11pt,a4paper]{article}
\usepackage[utf8]{inputenc}
\usepackage[T1]{fontenc}
\usepackage[ngerman]{babel}
\usepackage{geometry}
\usepackage{fancyhdr}
\usepackage{lastpage}
\usepackage{xcolor}
\usepackage{tcolorbox}
\usepackage{enumitem}
\usepackage{longtable}
\usepackage{hyperref}
\usepackage{amssymb}

\geometry{left=2.5cm, right=2.5cm, top=2.5cm, bottom=2.5cm}
\definecolor{fach}{HTML}{b35c00}
\definecolor{gray}{HTML}{666666}
\definecolor{successgreen}{HTML}{2a7a4b}
\definecolor{warnred}{HTML}{c62828}

\tcbuselibrary{skins,breakable}
\newtcolorbox{scorebox}{ colback=white, colframe=fach, fonttitle=\bfseries, title=Didaktik-Score, breakable }
\newtcolorbox{tippbox}[1][]{ colback=successgreen!5, colframe=successgreen, fonttitle=\bfseries, title=#1, breakable }

\pagestyle{fancy}
\fancyhf{}
\fancyhead[L]{\small\textcolor{gray}{Digitale Grundlagen -- 012 Suchmaschinen}}
\fancyhead[R]{\small\textcolor{gray}{LE Tier 2 / Profil S}}
\fancyfoot[C]{\small\thepage{} / \pageref{LastPage}}
\renewcommand{\headrulewidth}{0.4pt}

\begin{document}

\begin{titlepage}
    \centering
    \vspace*{2cm}
    {\Large\textcolor{gray}{Digitale Grundlagen -- Senioren 65+}}\\[2cm]
    {\Huge\textcolor{fach}{\textbf{Langentwurf}}}\\[0.5cm]
    {\large Tier 2 | Profil S}\\[2cm]
    {\LARGE\textbf{012 -- Suchmaschinen richtig nutzen}}\\[0.5cm]
    {\large Modul C: Internet \& Kommunikation}\\[2cm]
    \begin{tabular}{rl}
        \textbf{Zeitrahmen:} & 60--75 Minuten \\
        \textbf{Vorkenntnisse:} & Safari bedienen, Webseiten öffnen \\
    \end{tabular}
    \vfill
    {\normalsize Stand: \today}
\end{titlepage}

\tableofcontents
\newpage

\section{Bedingungsanalyse}

\subsection{Ausgangssituation}
\begin{itemize}
    \item Google ist oft die einzige bekannte Suchmaschine
    \item Suchbegriffe werden oft als ganze Sätze eingegeben
    \item Werbung wird nicht von echten Ergebnissen unterschieden
    \item Unsicherheit: ``Welchem Ergebnis kann ich trauen?''
    \item Überforderung durch zu viele Ergebnisse
\end{itemize}

\subsection{Typische Probleme}
\begin{itemize}
    \item ``Ich finde nicht, was ich suche''
    \item ``Da kommt nur Werbung''
    \item ``Woher weiß ich, ob das stimmt?''
    \item ``Warum zeigt Google mir andere Ergebnisse als meiner Freundin?''
\end{itemize}

\section{Sachanalyse}

\subsection{Kernkonzepte}
\begin{enumerate}
    \item \textbf{Suchmaschine:} Ein Programm, das das Internet nach deinen Begriffen durchsucht
    \item \textbf{Suchbegriffe (Keywords):} Die Wörter, nach denen gesucht wird
    \item \textbf{Suchergebnisse:} Die gefundenen Webseiten, sortiert nach Relevanz
    \item \textbf{Anzeigen (Werbung):} Bezahlte Ergebnisse -- gekennzeichnet mit ``Anzeige''
    \item \textbf{Vertrauenswürdigkeit:} Offizielle Seiten (.de, .gov) vs. unbekannte Quellen
\end{enumerate}

\subsection{Alltagsanalogien}
\begin{tabular}{|l|p{8cm}|}
\hline
\textbf{Begriff} & \textbf{Analogie} \\
\hline
Suchmaschine & Bibliothekarin, die für dich Bücher sucht \\
Suchbegriffe & Stichworte, die du der Bibliothekarin nennst \\
Suchergebnisse & Die Bücher, die sie dir bringt \\
Anzeigen & Bücher, für die der Verlag bezahlt hat (Reklame) \\
Filter & ``Nur Bücher aus diesem Jahr'' \\
\hline
\end{tabular}

\section{Didaktische Analyse}

\subsection{Stellung in der Reihe}
Aufbauend auf Browser-Grundlagen (011). Wichtig für alle weiteren Internet-Aktivitäten.

\subsection{Didaktische Reduktion}
\textbf{Ausgeklammert:} Erweiterte Suchoperatoren, Bildersuche, personalisierte Suche im Detail.

\textbf{Fokus:} Einfache, effektive Suche und Bewertung der Ergebnisse.

\section{Lernziele}

\begin{tabular}{|c|p{7cm}|c|c|}
\hline
\textbf{Nr.} & \textbf{Die Teilnehmer können...} & \textbf{Stufe} & \textbf{Niveau} \\
\hline
FZ1 & gute Suchbegriffe wählen (Schlüsselwörter statt Sätze) & Anwenden & Basis \\
FZ2 & Werbung von echten Ergebnissen unterscheiden & Verstehen & Basis \\
FZ3 & die Vertrauenswürdigkeit einer Webseite einschätzen & Anwenden & Basis \\
FZ4 & eine praktische Suche durchführen (z.B. Arzt finden) & Anwenden & Basis \\
FZ5 & alternative Suchmaschinen kennen (z.B. DuckDuckGo) & Wissen & Vertiefung \\
\hline
\end{tabular}

\section{Methodische Analyse}

\subsection{Suchstrategie für Anfänger}
\begin{tippbox}[So suchen Sie richtig]
\begin{enumerate}
    \item \textbf{Schlüsselwörter} statt ganzer Sätze: ``Hausarzt Rostock'' statt ``Wo finde ich einen guten Hausarzt in Rostock?''
    \item \textbf{Spezifisch} sein: ``Apfelkuchen Rezept einfach'' statt nur ``Kuchen''
    \item \textbf{Ergebnisse prüfen}: Wer hat das geschrieben? Ist es aktuell?
\end{enumerate}
\end{tippbox}

\section{Verlaufsplanung}

\begin{longtable}{|p{1cm}|p{2cm}|p{5cm}|p{3cm}|p{2cm}|}
\hline
\textbf{Zeit} & \textbf{Phase} & \textbf{Inhalt} & \textbf{TN-Aktivität} & \textbf{Material} \\
\hline
0--10 & Einstieg & ``Was suchen Sie oft im Internet?'' Sammeln & Berichten & -- \\
\hline
10--25 & Suchbegriffe & Gute vs. schlechte Suchbegriffe, Beispiele & Mitmachen & S.1 \\
\hline
25--40 & Ergebnisse lesen & Anzeigen erkennen, Titel und Beschreibung verstehen & Üben & S.1 \\
\hline
40--55 & Vertrauen & Welchen Seiten kann ich trauen? Checkliste & Diskutieren & S.2 \\
\hline
55--70 & Praxis & Eigene Suche: Arzt, Fahrplan, Rezept & Selbst suchen & S.2-3 \\
\hline
\end{longtable}

\section{Lösungen}

\textbf{Werbung erkennen:}
\begin{itemize}
    \item Steht ``Anzeige'' oder ``Ad'' dabei
    \item Oft ganz oben oder am Rand
    \item Leicht andersfarbiger Hintergrund
\end{itemize}

\textbf{Vertrauenswürdige Seiten:}
\begin{itemize}
    \item Offizielle Seiten: .de, .gov, .org
    \item Bekannte Institutionen: Verbraucherzentrale, Stiftung Warentest
    \item Zeitungen: tagesschau.de, zeit.de
    \item Vorsicht bei: Blogs, Foren, unbekannten Seiten
\end{itemize}

\textbf{Praktische Suchbeispiele:}
\begin{itemize}
    \item Arzt finden: ``Hausarzt [Ort] Bewertung''
    \item Fahrplan: ``Bahn [Start] [Ziel]'' oder direkt bahn.de
    \item Rezept: ``[Gericht] Rezept einfach''
    \item Öffnungszeiten: ``[Geschäft] [Ort] Öffnungszeiten''
\end{itemize}

\section{Didaktik-Score}

\begin{scorebox}
\begin{tabular}{|c|l|c|c|p{3.5cm}|}
\hline
\# & Dimension & Gew. & Score & Notiz \\
\hline
1 & Differenzierung & 15\% & 8/10 & Basis/Vertiefung vorhanden \\
2 & Taxonomie & 10\% & 9/10 & Verstehen + Anwenden \\
3 & Alltagsrelevanz & 10\% & 10/10 & Täglich gebraucht \\
4 & Aktivierung & 10\% & 10/10 & Eigene Suchen durchführen \\
5 & Scaffolding & 10\% & 9/10 & Suchstrategie-Box \\
6 & Feedback & 10\% & 8/10 & Beispiele mit Erklärung \\
7 & Sprachliche Klarheit & 10\% & 9/10 & Einfache Sprache \\
8 & Motivation & 10\% & 10/10 & Praktische Beispiele \\
9 & Barrierefreiheit & 10\% & 9/10 & Große Darstellung \\
10 & Praktikabilität & 5\% & 9/10 & 70 Min angemessen \\
\hline
\multicolumn{3}{|r|}{\textbf{GESAMT:}} & \textbf{9.15/10} & \textcolor{successgreen}{\textbf{FREIGABE}} \\
\hline
\end{tabular}
\end{scorebox}

\end{document}
